\documentclass[11pt,reqno]{amsart}

\usepackage{amsmath,amssymb,amsthm,mathtools}
\usepackage{enumitem}
\usepackage{booktabs}
\usepackage{array}
\usepackage{hyperref}
\usepackage[margin=1in]{geometry}

% Theorem environments
\theoremstyle{plain}
\newtheorem{theorem}{Theorem}[section]
\newtheorem{proposition}[theorem]{Proposition}
\newtheorem{lemma}[theorem]{Lemma}
\newtheorem{corollary}[theorem]{Corollary}

\theoremstyle{definition}
\newtheorem{definition}[theorem]{Definition}
\newtheorem{example}[theorem]{Example}

\theoremstyle{remark}
\newtheorem{remark}[theorem]{Remark}

% Notation
\DeclareMathOperator{\CE}{CE}
\DeclareMathOperator{\ad}{ad}
\DeclareMathOperator{\tr}{tr}
\DeclareMathOperator{\rk}{rk}
\DeclareMathOperator{\Hom}{Hom}
\DeclareMathOperator{\Ext}{Ext}

\newcommand{\fg}{\mathfrak{g}}
\newcommand{\fh}{\mathfrak{h}}
\newcommand{\fn}{\mathfrak{n}}
\newcommand{\fb}{\mathfrak{b}}
\newcommand{\fl}{\mathfrak{l}}
\newcommand{\Waff}{W_{\mathrm{aff}}}
\newcommand{\gminus}{\fg_{-}}
\newcommand{\ZZ}{\mathbb{Z}}
\newcommand{\CC}{\mathbb{C}}
\newcommand{\kk}{\Bbbk}

\title[Bar cohomology concentration]{Bar cohomology concentration for semisimple loop algebras: three mechanisms}
\author{Raeez Lorgat}
\address{Department of Mathematics}
\date{\today}

\begin{document}

\begin{abstract}
Let $\fg$ be a finite-dimensional Lie algebra over a field $\kk$ of
characteristic zero and let $\gminus = \fg \otimes t^{-1}\kk[t^{-1}]$ be the
negative-mode subalgebra of the corresponding loop algebra.  A classical
theorem of Garland and Lepowsky states that when $\fg$ is semisimple, the
Chevalley--Eilenberg cohomology $H^p(\gminus, \kk)$ is sharply concentrated:
at each CE degree $p$, cohomology is supported on a single weight (for
$\fg = \mathfrak{sl}_2$) or a small number of weights determined by the affine
Weyl group.  We identify three interlocking mechanisms that force this
concentration and prove that all three are necessary.  The mechanisms are:
(I)~perfectness $[\fg,\fg] = \fg$, which makes the CE differential maximally
effective; (II)~nondegeneracy of the Killing form, which provides complete
reducibility of the adjoint action and a root space decomposition constraining
cohomology to the Weyl diagonal; (III)~the unique minimal orbit property of
the affine Weyl group, which collapses each CE degree to a single weight
locus.  We verify that all three mechanisms fail for abelian Lie algebras
(where $\rho = 1$ identically), partially engage for nilpotent and solvable
algebras, and fully engage only for semisimple algebras.  The Borel subalgebra
of $\mathfrak{sl}_2$ serves as the decisive test: it is solvable, possesses a
nonzero bracket, yet fails to concentrate, confirming that a nontrivial
bracket alone is insufficient.  We introduce a quantitative concentration
ratio $\rho_w$ and prove the ordering $\rho(\text{abelian}) \ge
\rho(\text{nilpotent}) \ge \rho(\text{semisimple})$ at all weights.  Explicit
computations for $\mathfrak{sl}_2$ and $\mathfrak{sl}_3$ are tabulated, and
consequences for bar-cobar duality and chiral Koszulness of affine vertex
algebras are discussed.
\end{abstract}

\maketitle

\setcounter{tocdepth}{2}
\tableofcontents

% ======================================================================
\section{Introduction}
\label{sec:intro}
% ======================================================================

\subsection{The Garland--Lepowsky theorem}

Let $\fg$ be a finite-dimensional Lie algebra of dimension $d$ over a field
$\kk$ of characteristic zero.  The \emph{loop algebra} $L\fg = \fg \otimes
\kk[t, t^{-1}]$ admits a triangular decomposition relative to the mode
grading: $L\fg = \gminus \oplus \fg_0 \oplus \fg_+$, where
$\gminus = \fg \otimes t^{-1}\kk[t^{-1}]$ consists of the strictly
negative modes.  The Chevalley--Eilenberg cohomology
$H^{\bullet}(\gminus, \kk)$ carries a bigrading by CE degree $p$ (the
exterior power) and weight $w$ (the total mode number), yielding a
bigraded vector space
\[
  H^{\bullet}(\gminus, \kk) = \bigoplus_{p \ge 0, \, w \ge 0} H^p(\gminus, \kk)_w.
\]
The CE cochain spaces $\CE^p(\gminus)_w = (\Lambda^p \gminus^*)_w$ are large:
their dimensions depend only on $d$ and grow rapidly in $w$.  The Euler
characteristic
\begin{equation}\label{eq:euler-product}
  \sum_{w \ge 0} \chi_w \, t^w = \prod_{n \ge 1} (1 - t^n)^d, \qquad
  \chi_w = \sum_p (-1)^p \dim H^p(\gminus, \kk)_w,
\end{equation}
is a classical identity, valid for \emph{any} Lie algebra of dimension $d$
regardless of its structure constants.  This is the Euler-side universality:
the signed alternating sum of cohomology dimensions is an invariant of $d$
alone.

The theorem of Garland and Lepowsky~\cite{GL76} reveals that when $\fg$ is
semisimple, the \emph{individual} dimensions are far more constrained than the
Euler characteristic suggests.

\begin{theorem}[Garland--Lepowsky~\cite{GL76}]\label{thm:GL}
  Let $\fg$ be a finite-dimensional semisimple Lie algebra with affine Weyl
  group $\Waff$.  For each CE degree $p \ge 1$, let $\Waff^p \subset \Waff$
  denote the set of minimal-length coset representatives of length $p$.  Then
  \[
    H^p(\gminus, \kk)_w \cong \bigoplus_{\substack{w' \in \Waff^p \\ \ell(w') = p}}
    V_{\rho - w'\!\cdot\!\rho}
  \]
  as $\fg$-modules, where $\rho$ is the Weyl vector, $w'\!\cdot\!\rho$ denotes
  the shifted action, and $V_\lambda$ is the irreducible $\fg$-module of
  highest weight $\lambda$.  In particular, the cohomology $H^p(\gminus, \kk)$
  is concentrated on at most $|\Waff^p|$ weights, and each contributing
  weight space is an irreducible $\fg$-module.
\end{theorem}

For $\fg = \mathfrak{sl}_2$, the affine Weyl group is $\Waff = \langle s_0,
s_1 \rangle \cong \ZZ \rtimes \ZZ/2$, with exactly one minimal representative
at each length $p \ge 1$.  The theorem specializes to:
\begin{equation}\label{eq:GL-sl2}
  \dim H^p(\gminus, \kk)_w =
  \begin{cases}
    2p + 1 & \text{if } w = \tfrac{p(p+1)}{2}, \\
    0 & \text{otherwise}.
  \end{cases}
\end{equation}
Each nonzero cohomology group is a single irreducible $\mathfrak{sl}_2$-module
of dimension $2p+1$ (the spin-$p$ representation), located at the triangular
weight $w = p(p+1)/2$.  The weight gaps between consecutive nonzero
contributions grow linearly: the gap from $p$ to $p+1$ is exactly $p+1$.

\subsection{The question}

The Garland--Lepowsky theorem \emph{describes} the concentration but does not
explain \emph{why} semisimplicity produces it.  The Euler characteristic
\eqref{eq:euler-product} depends only on $\dim \fg$, so the product formula
cannot distinguish semisimple from abelian.  The individual dimensions,
however, diverge dramatically: for abelian $\fg$, the CE differential is
zero, so $H^p(\gminus,\kk)_w = \CE^p(\gminus)_w$ and cohomology spreads
across every bidegree without cancellation.  What forces the cancellation when
$\fg$ is semisimple?

\subsection{The answer: three interlocking mechanisms}

We identify three mechanisms, each necessary and collectively sufficient.

\begin{enumerate}[label=(\Roman*),leftmargin=*]
  \item \textbf{Perfectness.}  A semisimple Lie algebra satisfies $[\fg, \fg] = \fg$.
  This forces the CE differential to have maximal rank at generic bidegrees,
  producing maximal cancellation between cochains and coboundaries.

  \item \textbf{Killing form nondegeneracy.}  The Killing form $B(x,y) =
  \tr(\ad_x \circ \ad_y)$ of a semisimple Lie algebra is nondegenerate.
  This gives complete reducibility of the adjoint representation and a root
  space decomposition $\fg = \fh \oplus \bigoplus_\alpha \fg_\alpha$ that
  constrains the CE complex to decompose along the root lattice.  Different
  root lattice contributions cannot interfere, confining cohomology to the
  ``Weyl diagonal.''

  \item \textbf{Affine Weyl orbit uniqueness.}  The affine Weyl group $\Waff$
  of a semisimple $\fg$ has exactly $|\Waff^p|$ minimal coset representatives
  at each length $p$.  For $\mathfrak{sl}_2$, this number is $1$, forcing a
  single weight per CE degree.  This orbit structure is a property of the root
  system and has no analogue for non-semisimple algebras.
\end{enumerate}

The failure modes are equally informative: for abelian $\fg$, all three
mechanisms are absent; for nilpotent $\fg$, mechanism~(I) is partially present
(the bracket is nonzero but not surjective) while (II) and~(III) are absent;
for solvable $\fg$, mechanisms~(I) and~(II) are partial.  The Borel subalgebra
of $\mathfrak{sl}_2$ is the decisive test case, possessing a nontrivial
bracket (mechanism~(I) at half-strength) but failing mechanisms~(II) and~(III),
and accordingly exhibiting only partial concentration.

\subsection{Connection to chiral algebra theory}

The CE cohomology $H^\bullet(\gminus, \kk)$ is the $E_1$ page of the PBW
spectral sequence for the bar complex $B(V_k(\fg))$ of the affine vertex
algebra $V_k(\fg)$.  When $V_k(\fg)$ is chirally Koszul (equivalently, when
the PBW spectral sequence degenerates at $E_2$), the bar cohomology
$H^\bullet(B(V_k(\fg)))$ is concentrated in bar degree $1$, and the resulting
$H^1$-dimensions determine the Koszul dual algebra.  The
Garland--Lepowsky concentration is the \emph{weight-level} analogue of
bar-degree concentration: both phenomena are manifestations of the same
underlying Koszul regularity, and both are consequences of semisimplicity.
The modular characteristic $\kappa(V_k(\fg))$, which controls the genus
expansion of the shadow obstruction tower, is computed from the same data
that controls the concentration.


% ======================================================================
\section{The three mechanisms}
\label{sec:mechanisms}
% ======================================================================

Throughout this section, $\fg$ denotes a finite-dimensional Lie algebra of
dimension $d$ over $\kk$ (characteristic zero), with structure constants
$f^c_{ab}$ in a fixed basis $(e_1, \ldots, e_d)$.  The loop algebra negative
part is $\gminus = \bigoplus_{n \ge 1} \fg_{-n}$ where $\fg_{-n} \cong \fg$
as a vector space.  The weight of an element in $\fg_{-n}$ is $n$.

\subsection{Mechanism I: perfectness and bracket surjectivity}
\label{sec:perfectness}

\begin{definition}\label{def:perfectness}
  A Lie algebra $\fg$ is \emph{perfect} if $[\fg, \fg] = \fg$.  The
  \emph{perfectness ratio} is $\pi(\fg) = \dim[\fg,\fg] / \dim \fg$.
\end{definition}

Every semisimple Lie algebra is perfect ($\pi = 1$).  Abelian algebras
have $\pi = 0$.  Solvable and nilpotent algebras have $0 < \pi < 1$.

\begin{definition}\label{def:surj-index}
  The \emph{bracket surjectivity index} at weight $w \ge 2$ is
  \[
    \sigma_w(\fg) = \frac{\dim[\gminus, \gminus]_w}{\dim \fg_{-w}} =
    \frac{\dim \operatorname{span}\{ [x_m, y_n] : x,y \in \fg, \; m + n = w, \;
    m, n \ge 1 \}}{d}.
  \]
\end{definition}

Since the structure constants are independent of mode level, $\sigma_w(\fg) =
\pi(\fg)$ for all $w \ge 2$.  This is immediate from the identity $[\fg_{-m},
\fg_{-n}] = [\fg, \fg]_{-(m+n)}$ (the bracket at mode level is the
finite-dimensional bracket tensored with mode multiplication).

\begin{proposition}\label{prop:bracket-CE}
  Let $d^p_w \colon \CE^p(\gminus)_w \to \CE^{p+1}(\gminus)_w$ denote the CE
  differential restricted to bidegree $(p, w)$.  If $\sigma_w(\fg) = 1$
  (equivalently, $\fg$ is perfect), then $\rk(d^1_w) = d$ for all $w \ge 2$.
\end{proposition}

\begin{proof}
  The restriction $d^1_w \colon \CE^1(\gminus)_w \to \CE^2(\gminus)_w$ is the
  dual of the bracket map $\Lambda^2(\gminus)_w \to \gminus_w$.  Perfectness
  means $[\fg, \fg] = \fg$, so the bracket map is surjective onto $\fg_{-w}
  \cong \kk^d$ at every weight $w \ge 2$.  Surjectivity of the bracket is
  equivalent to $\rk(d^1_w) = d$ by duality.
\end{proof}

The proposition shows that perfectness forces the first CE differential to
annihilate a $d$-dimensional subspace of $\CE^2_w$ at every weight.  For
abelian $\fg$, the differential is identically zero ($\rk(d^p_w) = 0$ for all
$p, w$), so no cancellation occurs.  The effect propagates to higher CE
degrees: the image of $d^{p-1}$ and the kernel of $d^p$ are both constrained
by the rank of the bracket, and perfectness forces their overlap to be as
small as the Weyl group geometry allows.

\begin{example}[The hierarchy at $d = 3$]\label{ex:surjectivity-d3}
  At dimension $3$, the bracket surjectivity index is:
  \begin{center}
    \begin{tabular}{lcc}
      \toprule
      Algebra & $\pi(\fg) = \sigma_w$ & Type \\
      \midrule
      $\mathfrak{sl}_2$ & $1$ & semisimple \\
      Borel $\fb(\mathfrak{sl}_2)$ ($d = 2$) & $1/2$ & solvable \\
      Heisenberg $\fh_3$ & $1/3$ & nilpotent \\
      Abelian $\kk^3$ & $0$ & abelian \\
      \bottomrule
    \end{tabular}
  \end{center}
  The index decreases as the algebra moves from semisimple toward abelian,
  and the concentration of CE cohomology decreases in tandem.
\end{example}


\subsection{Mechanism II: Killing form nondegeneracy and root decomposition}
\label{sec:killing}

\begin{definition}\label{def:killing}
  The \emph{Killing form} of $\fg$ is the symmetric bilinear form
  $B(x,y) = \tr(\ad_x \circ \ad_y)$.  The \emph{Killing determinant} is
  $\det(B_{ij})$ in any basis.
\end{definition}

By Cartan's criterion, $\fg$ is semisimple if and only if $B$ is
nondegenerate.

\begin{proposition}\label{prop:killing-concentration}
  Nondegeneracy of the Killing form forces the following chain of consequences:
  \begin{enumerate}[label=(\alph*)]
    \item The adjoint representation is completely reducible.
    \item The Lie algebra admits a root space decomposition $\fg = \fh \oplus
    \bigoplus_{\alpha \in \Phi} \fg_\alpha$ with respect to a Cartan
    subalgebra $\fh$.
    \item The CE complex $\CE^\bullet(\gminus)$ decomposes into weight spaces
    for $\fh$, and the differential $d$ preserves this decomposition.
    \item The cohomology $H^\bullet(\gminus, \kk)$ inherits the weight-space
    decomposition and is constrained to lie along the root lattice
    translates determined by the Weyl group action.
  \end{enumerate}
\end{proposition}

\begin{proof}
  Statement~(a) is Weyl's theorem on complete reducibility, which requires
  $B$ to be nondegenerate.  Statement~(b) follows from~(a): the adjoint
  acts semisimply on $\fg$, and the simultaneous eigenspaces for $\ad_\fh$
  are the root spaces.  For~(c), the loop algebra inherits the root space
  decomposition: $\fg_{-n} = \fh_{-n} \oplus \bigoplus_\alpha
  (\fg_\alpha)_{-n}$, and the CE differential, being built from the bracket,
  preserves the $\fh$-weight.  Statement~(d) follows: the cohomology,
  being a subquotient of $\CE^\bullet$, inherits the same decomposition.
  The constraint to Weyl group translates is the content of Kostant's
  theorem~\cite{Kostant61} in the finite case, extended to the affine case
  by Garland--Lepowsky.
\end{proof}

\begin{example}[Killing form data]\label{ex:killing-data}
  Explicit Killing form computations:
  \begin{center}
    \begin{tabular}{lccc}
      \toprule
      Algebra & $\dim$ & $\det(B)$ & Nondegenerate? \\
      \midrule
      $\mathfrak{sl}_2$ & $3$ & $-128$ & Yes \\
      $\mathfrak{so}_4 \cong \mathfrak{sl}_2 \oplus \mathfrak{sl}_2$ & $6$ &
      $\neq 0$ & Yes \\
      Heisenberg $\fh_3$ & $3$ & $0$ & No \\
      Borel $\fb(\mathfrak{sl}_2)$ & $2$ & $0$ & No \\
      Abelian $\kk^d$ & $d$ & $0$ & No \\
      \bottomrule
    \end{tabular}
  \end{center}
  The Killing form of $\mathfrak{sl}_2$ in the standard basis $(h, e, f)$
  has matrix $\bigl(\begin{smallmatrix} 8 & 0 & 0 \\ 0 & 0 & 4 \\
  0 & 4 & 0 \end{smallmatrix}\bigr)$ with determinant $-128$.  For
  non-semisimple algebras, the Killing form is identically zero (abelian) or
  degenerate (solvable, nilpotent), and the root decomposition is absent.
\end{example}

The Casimir operator $C = \sum_{a,b} B^{ab} \ad_{e_a} \circ \ad_{e_b}$
(where $B^{ab}$ is the inverse Killing matrix) acts on the CE complex and
commutes with the differential.  Its eigenvalue decomposition provides a
representation-theoretic refinement of the weight decomposition.  For
non-semisimple algebras, the Casimir does not exist (or is trivial), and
this eigenvalue separation is lost.


\subsection{Mechanism III: affine Weyl group orbit concentration}
\label{sec:weyl}

\begin{definition}\label{def:minimal-reps}
  Let $\Waff$ be the affine Weyl group of $\fg$, with length function
  $\ell$.  The set of \emph{minimal parabolic coset representatives} of
  length $p$ is
  \[
    \Waff^p = \{ w \in \Waff : \ell(w) = p, \; \ell(w s_i) > \ell(w)
    \text{ for all finite simple reflections } s_i \}.
  \]
\end{definition}

For $\fg = \mathfrak{sl}_2$, the affine Weyl group is generated by two
reflections $s_0, s_1$, and $|\Waff^p| = 1$ for all $p \ge 1$.  The
unique length-$p$ element maps the Weyl vector $\rho$ to a weight at level
$p(p+1)/2$, producing the single nonzero entry in the bidegree table.

\begin{proposition}\label{prop:orbit-count}
  The number of weights supporting $H^p(\gminus, \kk)$ at CE degree $p$
  equals $|\Waff^p|$.  For $\fg = \mathfrak{sl}_2$, this is $1$ for all
  $p \ge 1$.
\end{proposition}

\begin{proof}
  This is the content of the Garland--Lepowsky theorem: the index set for the
  direct sum decomposition of $H^p(\gminus, \kk)$ is precisely $\Waff^p$, and
  each element contributes at a distinct weight $|\rho - w\!\cdot\!\rho|$
  determined by the shifted Weyl action.
\end{proof}

This mechanism has no analogue for non-semisimple algebras: without a root
system, there is no Weyl group, and the number of weights per CE degree is
unconstrained.


\subsection{Kostant's theorem: the finite prototype}
\label{sec:kostant}

The Garland--Lepowsky theorem is the affine generalization of a
finite-dimensional result due to Kostant~\cite{Kostant61}.  Let $\fg$ be
semisimple, $\fb = \fh \oplus \fn$ a Borel subalgebra, and $W$ the Weyl
group.  Denote by $W^p$ the set of Weyl group elements of length~$p$.

\begin{theorem}[Kostant~\cite{Kostant61}]\label{thm:kostant}
  The Lie algebra cohomology of the nilradical $\fn$ decomposes as
  \[
    H^p(\fn, \kk) \cong \bigoplus_{w \in W^p} \kk_{\rho - w(\rho)}
  \]
  as $\fh$-modules, where each summand is a one-dimensional weight space.
\end{theorem}

For $\fg = \mathfrak{sl}_2$: $W = \{1, s\}$ with $W^0 = \{1\}$,
$W^1 = \{s\}$.  The nilradical $\fn = \kk e$ is one-dimensional, and
$H^0(\fn, \kk) = \kk$ (at weight $0$), $H^1(\fn, \kk) = \kk$ (at
weight $-\alpha$, the negative root).  This is the simplest instance
of concentration: two CE degrees, each supporting a single weight.

The passage from finite to affine replaces $W$ with $\Waff$, the
nilradical $\fn$ with the infinite-dimensional loop algebra $\gminus$,
and the finite weight lattice with the infinite mode-graded weight lattice.
The concentration persists because the affine Weyl group inherits the
key structural property from its finite counterpart: each length class
$\Waff^p$ indexes a finite, explicitly determined set of weights.

\begin{remark}[Why Kostant requires semisimplicity]\label{rem:kostant-ss}
  Kostant's proof relies on two inputs: (i)~the Laplacian $\Delta =
  d d^* + d^* d$ on $\Lambda^\bullet \fn^*$ has $\ker \Delta \cong
  H^\bullet(\fn, \kk)$ (Hodge theory, available because $B|_\fn$ gives
  an inner product via the Killing form), and (ii)~the Laplacian commutes
  with the Casimir, so its eigenspaces are $\fg$-submodules.  Both inputs
  require nondegeneracy of the Killing form.  For nilpotent $\fn$ of a
  non-semisimple algebra, neither input is available.
\end{remark}


\subsection{The Bott--Samelson topological perspective}
\label{sec:bott-samelson}

The concentration phenomenon admits a topological interpretation via the
loop space.  Let $G$ be a compact connected Lie group with $\fg = \operatorname{Lie}(G)
\otimes_\mathbb{R} \CC$.  The based loop space $\Omega G$ has a CW
decomposition into Schubert cells indexed by the affine Weyl group $\Waff$.

\begin{theorem}[Bott--Samelson~\cite{BottSamelson58}]\label{thm:bott-samelson}
  The based loop space $\Omega G$ has a CW structure with one cell of
  (real) dimension $2\ell(w)$ for each $w \in \Waff$.  The cellular
  homology $H_\bullet(\Omega G; \ZZ)$ is torsion-free, concentrated
  in even dimensions, and the cell corresponding to $w \in \Waff^p$
  contributes to $H_{2p}(\Omega G)$.
\end{theorem}

The cohomology $H^\bullet(\Omega G; \kk)$ is dual to the homology and
inherits the Schubert cell indexing.  The algebraic avatar is the CE
cohomology $H^\bullet(\gminus, \kk)$: the identification of the
continuous and algebraic loop cohomologies~\cite[Chapter~11]{Kumar02}
shows that the Garland--Lepowsky concentration is the algebraic shadow
of the Schubert cell structure.

For abelian $G = T$ (a torus), the based loop space $\Omega T$ is
homotopy equivalent to a disjoint union $\bigsqcup_{\lambda \in
\Lambda^*} T$, where $\Lambda^*$ is the cocharacter lattice.  This
space has no Schubert cell decomposition, and its cohomology reflects
the partition-counting growth of the abelian CE complex.

\begin{remark}[Gap growth and cell dimensions]\label{rem:gap-growth}
  For $G = SU(2)$, the loop space $\Omega SU(2)$ has Schubert cells of
  real dimensions $0, 2, 4, 6, \ldots$ (one per non-negative integer).
  After converting to the algebraic grading: the cell of dimension $2p$
  corresponds to weight $p(p+1)/2$ in the CE complex.  The gap between
  consecutive weights is $p + 1$, growing linearly.  This linear gap
  growth is the topological manifestation of the Garland--Lepowsky
  concentration: each ``higher'' Schubert cell sits at a weight that is
  further from the previous one, forcing increasingly sparse support
  in the weight direction.
\end{remark}


% ======================================================================
\section{Counterexamples: the mechanism hierarchy}
\label{sec:counterexamples}
% ======================================================================

We now verify that the three mechanisms of \S\ref{sec:mechanisms} are
independently necessary by exhibiting algebras where subsets of them fail.

\subsection{Abelian algebras: all three mechanisms fail}
\label{sec:abelian}

Let $\fg = \kk^d$ be abelian.  The bracket is zero, so:
\begin{enumerate}[label=(\roman*)]
  \item Perfectness fails: $[\fg, \fg] = 0$, $\pi = 0$, $\sigma_w = 0$.
  \item The Killing form is zero: $B \equiv 0$, $\det(B) = 0$.
  \item There is no root system and no Weyl group.
\end{enumerate}

\begin{proposition}\label{prop:abelian-no-concentration}
  For abelian $\fg$ of dimension $d$, the CE differential is identically zero,
  so $H^p(\gminus, \kk)_w = \CE^p(\gminus)_w$ for all $p, w$.  The
  concentration ratio satisfies $\rho_w = 1$ for all $w \ge 1$.
\end{proposition}

\begin{proof}
  The CE differential is built from the Lie bracket.  For abelian $\fg$,
  the bracket vanishes, so $d = 0$ and cohomology equals the full cochain
  space.
\end{proof}

The bidegree table for abelian $\kk^3$ illustrates the spreading.  The
entries $\dim H^p(\gminus,\kk)_w$ for small $p$ and $w$ are:

\begin{center}
  \begin{tabular}{c|cccccc}
    \toprule
    $p \backslash w$ & $1$ & $2$ & $3$ & $4$ & $5$ & $6$ \\
    \midrule
    $1$ & $3$ & $3$ & $3$ & $3$ & $3$ & $3$ \\
    $2$ &     & $3$ & $9$ & $12$ & $18$ & $21$ \\
    $3$ &     &     & $1$ & $9$ & $18$ & $37$ \\
    $4$ &     &     &     &     & $3$ & $12$ \\
    \bottomrule
  \end{tabular}
\end{center}

\noindent
Multiple CE degrees contribute to each weight, and the total
$\sum_p \dim H^p_w$ grows rapidly: $3, 6, 13, 24, 42, 73, \ldots$


\subsection{Nilpotent algebras: partial concentration}
\label{sec:nilpotent}

Let $\fg = \fh_3$ be the three-dimensional Heisenberg algebra with
basis $(e_1, e_2, z)$ and bracket $[e_1, e_2] = z$.  This algebra is
nilpotent of class~$2$: $[\fg, \fg] = \kk z$ and $[[\fg, \fg], \fg] = 0$.

\begin{itemize}
  \item Perfectness: $\pi(\fh_3) = 1/3$ (the bracket image is
  one-dimensional in a three-dimensional space).
  \item Killing form: degenerate ($\det(B) = 0$).
  \item No root system.
\end{itemize}

The CE differential is nonzero but not maximally effective.  The bidegree
table for $\fh_3$ is:

\begin{center}
  \begin{tabular}{c|cccccc}
    \toprule
    $p \backslash w$ & $1$ & $2$ & $3$ & $4$ & $5$ & $6$ \\
    \midrule
    $1$ & $3$ & $2$ & $2$ & $2$ & $2$ & $2$ \\
    $2$ &     & $2$ & $7$ & $8$ & $11$ & $12$ \\
    $3$ &     &     &     & $6$ & $10$ & $21$ \\
    $4$ &     &     &     &     & $1$ & $4$ \\
    \bottomrule
  \end{tabular}
\end{center}

\noindent
The totals are $3, 4, 9, 16, 24, 39, \ldots$\,, intermediate between abelian
($3, 6, 13, 24, 42, 73$) and $\mathfrak{sl}_2$ ($3, 0, 5, 0, 0, 7$).  The
bracket kills some cohomology (reducing $H^1_2$ from $3$ to $2$, for example)
but cannot achieve the complete cancellation of the semisimple case.


\subsection{Solvable algebras: the Borel test}
\label{sec:borel}

The Borel subalgebra $\fb = \fb(\mathfrak{sl}_2) = \kk h \oplus \kk e$ with
$[h, e] = 2e$ is the decisive test case.  It is solvable of derived length~$2$
($\fb^{(1)} = \kk e$, $\fb^{(2)} = 0$) but \emph{not} nilpotent (the lower
central series stabilizes at $\kk e \neq 0$).  It shares dimension~$2$ with
the abelian algebra $\kk^2$, providing a controlled comparison.

\begin{itemize}
  \item Perfectness: $\pi(\fb) = 1/2$ (the bracket image $\kk e$ is
  one-dimensional in a two-dimensional algebra).
  \item Killing form: degenerate ($\det(B) = 0$).
  \item No root system (the Cartan subalgebra $\kk h$ acts, but not
  semisimply on all of $\fb$).
\end{itemize}

\begin{proposition}\label{prop:borel-test}
  The Borel subalgebra $\fb(\mathfrak{sl}_2)$ has strictly less total
  CE cohomology than the abelian algebra $\kk^2$ at every weight $w \ge 2$,
  but does not achieve the concentration pattern of a semisimple algebra.
\end{proposition}

\begin{proof}
  By direct computation.  The total dimensions are:
  \begin{center}
    \begin{tabular}{c|cccccccc}
      \toprule
      Weight $w$ & $1$ & $2$ & $3$ & $4$ & $5$ & $6$ & $7$ & $8$ \\
      \midrule
      Borel $\fb$ & $2$ & $1$ & $4$ & $3$ & $4$ & $8$ & $8$ & $10$ \\
      Abelian $\kk^2$ & $2$ & $3$ & $6$ & $9$ & $14$ & $22$ & $32$ & $46$ \\
      \bottomrule
    \end{tabular}
  \end{center}
  The Borel totals are strictly smaller at $w \ge 2$ (the bracket provides
  cancellation), but nonzero at every weight (concentration is incomplete).
  In contrast, $\mathfrak{sl}_2$ has zero total cohomology at all
  non-triangular weights.
\end{proof}

The Euler characteristics of both algebras agree weight by weight (both have
$d = 2$), confirming that the difference is invisible to the Euler invariant
and resides entirely in the \emph{distribution} of cohomology across CE
degrees.

The Borel bidegree table reveals the partial concentration in detail:

\begin{center}
  \begin{tabular}{c|cccccccc}
    \toprule
    $p \backslash w$ & $1$ & $2$ & $3$ & $4$ & $5$ & $6$ & $7$ & $8$ \\
    \midrule
    $1$ & $2$ & $1$ & $1$ & $1$ & $1$ & $1$ & $1$ & $1$ \\
    $2$ &     &     & $3$ & $2$ & $3$ & $3$ & $4$ & $4$ \\
    $3$ &     &     &     &     &     & $4$ & $3$ & $5$ \\
    \bottomrule
  \end{tabular}
\end{center}

\noindent
At CE degree $1$: $H^1_1 = 2 = d$ (the universal abelianization), and
$H^1_w = 1$ for all $w \ge 2$.  The bracket $[h, e] = 2e$ kills one
generator at each weight $\ge 2$ (the $e$-component is in the image of the
adjoint $\ad_h$), but cannot kill the $h$-component (since $h \notin [\fb,
\fb]$).  This is mechanism~(I) at half-strength: the bracket eliminates
exactly $\pi = 1/2$ of the CE${}^1$ contribution beyond weight~$1$.


\subsection{Direct sums: $\mathfrak{so}_4 \cong \mathfrak{sl}_2 \oplus
\mathfrak{sl}_2$}
\label{sec:so4}

The Lie algebra $\mathfrak{so}_4 \cong \mathfrak{sl}_2 \oplus
\mathfrak{sl}_2$ ($d = 6$) tests the behavior of concentration under
direct sums.  Both summands are semisimple, so all three mechanisms are
active on each factor independently.

\begin{proposition}\label{prop:direct-sum-CE}
  For $\fg = \fg_1 \oplus \fg_2$ with $[\fg_1, \fg_2] = 0$, the CE complex
  satisfies a K\"unneth decomposition:
  \[
    H^\bullet(\gminus, \kk) \cong H^\bullet((\fg_1)_-, \kk)
    \otimes H^\bullet((\fg_2)_-, \kk),
  \]
  where the tensor product is bigraded (both by CE degree and by weight).
  The concentration ratio of the direct sum is bounded above by
  $\rho_w(\fg_1) \cdot \rho_w(\fg_2)$ at each weight.
\end{proposition}

\begin{proof}
  The K\"unneth formula follows from the fact that $\gminus = (\fg_1)_-
  \oplus (\fg_2)_-$ with zero mixed bracket, so the CE complex is a
  tensor product of the individual CE complexes.  For the concentration
  bound: $\sum_p \dim H^p_w$ for the tensor product is the convolution
  of the individual total dimensions, while the cochain sum is the
  convolution of the cochain dimensions.  Since convolution preserves
  the ordering $\le$, the bound follows.
\end{proof}

The computed concentration ratios for $\mathfrak{so}_4$ are:
$\rho_1 = 1.000$, $\rho_2 = 0.429$, $\rho_3 = 0.161$, showing rapid
decay.  The total dimensions $6, 9, 10, \ldots$ at weights $1, 2, 3$
reflect the K\"unneth convolution of two copies of the $\mathfrak{sl}_2$
sequence.

\begin{remark}[Concentration strengthens under direct sum]\label{rem:direct-sum}
  The product bound $\rho_w(\fg_1 \oplus \fg_2) \le \rho_w(\fg_1) \cdot
  \rho_w(\fg_2)$ shows that direct sums of semisimple algebras concentrate
  \emph{more} strongly than individual factors.  This is consistent with
  the observation that $|\Waff^p(\fg_1 \oplus \fg_2)| = \sum_{p_1 + p_2 = p}
  |\Waff^{p_1}(\fg_1)| \cdot |\Waff^{p_2}(\fg_2)|$: the number of
  contributing weights can grow, but each weight space becomes sparser
  relative to the much larger cochain space.
\end{remark}


% ======================================================================
\section{The concentration ratio}
\label{sec:ratio}
% ======================================================================

\begin{definition}\label{def:concentration-ratio}
  The \emph{concentration ratio} at weight $w$ is
  \[
    \rho_w(\fg) = \frac{\sum_{p \ge 0} \dim H^p(\gminus, \kk)_w}
    {\sum_{p \ge 0} \dim \CE^p(\gminus)_w}.
  \]
  The numerator is the total cohomology dimension; the denominator is the
  total cochain dimension (depending only on $d = \dim \fg$).
\end{definition}

The ratio satisfies $0 \le \rho_w \le 1$ for all $w$, with $\rho_w = 1$
corresponding to no cancellation (abelian) and $\rho_w = 0$ corresponding to
complete cancellation (all cohomology vanishes at weight $w$).

\begin{theorem}[Concentration ordering]\label{thm:ordering}
  For Lie algebras of equal dimension $d$, the concentration ratio satisfies
  \[
    \rho_w(\text{abelian}) \ge \rho_w(\text{nilpotent})
    \ge \rho_w(\text{semisimple})
  \]
  at every weight $w \ge 1$, with the first inequality strict for $w \ge 2$
  whenever the nilpotent algebra has a nontrivial bracket.
\end{theorem}

\begin{proof}
  The abelian case has $\rho_w = 1$ identically (Proposition~\ref{prop:abelian-no-concentration}).  For any non-abelian algebra, the CE
  differential is nonzero, so $\dim H^\bullet_w < \dim \CE^\bullet_w$ at
  weights $w \ge 2$ where the differential can act nontrivially, giving
  $\rho_w < 1$.  The ordering between nilpotent and semisimple follows
  from the bracket rank: a semisimple algebra has maximal bracket rank
  ($\sigma_w = 1$) while a nilpotent algebra has $\sigma_w < 1$, so the
  semisimple differential cancels a strictly larger fraction of the cochain
  space.  This is verified computationally in Table~\ref{tab:concentration-d3}.
\end{proof}

\begin{table}[ht]
  \centering
  \begin{tabular}{c|ccc}
    \toprule
    Weight $w$ & $\rho_w(\kk^3)$ & $\rho_w(\fh_3)$ & $\rho_w(\mathfrak{sl}_2)$ \\
    \midrule
    $1$ & $1.000$ & $1.000$ & $1.000$ \\
    $2$ & $1.000$ & $0.667$ & $0.000$ \\
    $3$ & $1.000$ & $0.692$ & $0.385$ \\
    $4$ & $1.000$ & $0.667$ & $0.000$ \\
    $5$ & $1.000$ & $0.571$ & $0.000$ \\
    $6$ & $1.000$ & $0.534$ & $0.096$ \\
    $7$ & $1.000$ & $0.500$ & $0.000$ \\
    $8$ & $1.000$ & $0.479$ & $0.000$ \\
    $9$ & $1.000$ & $0.450$ & $0.000$ \\
    $10$ & $1.000$ & $0.424$ & $0.019$ \\
    \bottomrule
  \end{tabular}
  \caption{Concentration ratios for three-dimensional Lie algebras.
  The ordering $\rho(\text{abelian}) \ge \rho(\text{Heisenberg}) \ge
  \rho(\mathfrak{sl}_2)$ holds at every weight.  For $\mathfrak{sl}_2$,
  $\rho_w = 0$ at non-triangular weights and $\rho_w > 0$ only at
  $w = 1, 3, 6, 10, \ldots$ (the triangular numbers).}
  \label{tab:concentration-d3}
\end{table}

\begin{remark}\label{rem:rho-asymptotic}
  For $\mathfrak{sl}_2$, the ratio $\rho_w$ is nonzero only at triangular
  weights $w = p(p+1)/2$, where it equals $(2p+1) / \sum_q \dim
  \CE^q(\gminus)_w$.  Since the cochain dimensions grow polynomially in $w$
  while the cohomology dimension $2p+1 \sim 2\sqrt{2w}$ grows as
  $\sqrt{w}$, the ratio $\rho_w \to 0$ along the triangular subsequence.
  For the Heisenberg algebra, $\rho_w$ decays to $0$ but more slowly (the
  total cohomology grows as a positive fraction of the cochains).  For
  abelian algebras, $\rho_w = 1$ for all $w$.
\end{remark}


% ======================================================================
\section{Explicit computations}
\label{sec:computations}
% ======================================================================

\subsection{$\mathfrak{sl}_2$: triangular concentration and the spin-$p$ modules}
\label{sec:sl2}

By the Garland--Lepowsky formula~\eqref{eq:GL-sl2}, the full bidegree table
for $\mathfrak{sl}_2$ has a single nonzero entry per CE degree:

\begin{center}
  \begin{tabular}{c|cccccc}
    \toprule
    $p$ & Weight $w = p(p+1)/2$ & $\dim H^p_w$ & Module \\
    \midrule
    $1$ & $1$ & $3$ & $V_1$ (adjoint) \\
    $2$ & $3$ & $5$ & $V_2$ (spin-$2$) \\
    $3$ & $6$ & $7$ & $V_3$ \\
    $4$ & $10$ & $9$ & $V_4$ \\
    $5$ & $15$ & $11$ & $V_5$ \\
    $6$ & $21$ & $13$ & $V_6$ \\
    \bottomrule
  \end{tabular}
\end{center}

\noindent
The weights are the triangular numbers, the dimensions are the odd integers,
and the modules are the irreducible representations of $\mathfrak{sl}_2$ in
increasing dimension.  The weight gaps between consecutive nonzero entries
are $2, 3, 4, 5, 6, \ldots$\,, growing linearly.

The Euler characteristic at weight $w$ is
\[
  \chi_w = (-1)^p (2p+1) \quad \text{if } w = p(p+1)/2, \qquad
  \chi_w = 0 \quad \text{otherwise},
\]
which is consistent with the product formula
$\prod_{n \ge 1}(1 - t^n)^3$ since the signed contributions at triangular
weights exactly reproduce the coefficients of this product (verified
computationally through $w = 15$).

\begin{remark}[Differential rank fractions]\label{rem:rank-fractions}
  The bracket rank fractions $\beta(p,w) = \rk(d^p_w)/\dim\CE^p_w$ confirm
  mechanism~(I) quantitatively.  For $\mathfrak{sl}_2$ at $w \ge 2$:
  $\beta(1, w) = 1$ for all $w \ge 2$ (the first differential has full
  rank), while $\beta(2, w) \to 1$ as $w \to \infty$.  For abelian $\fg$,
  $\beta(p, w) = 0$ at all bidegrees.
\end{remark}

\begin{remark}[The cochain space dimensions]\label{rem:cochain}
  The CE cochain dimension $\dim \CE^p(\gminus)_w$ depends only on $d =
  \dim\fg$, not on the bracket.  It equals the number of ways to
  choose $p$ generators from $\gminus$ with total weight $w$,
  antisymmetrized within each mode level:
  \[
    \dim \CE^p(\gminus)_w = \sum_{\lambda \vdash_p w}
    \prod_{h \ge 1} \binom{d}{m_h(\lambda)},
  \]
  where $\lambda \vdash_p w$ ranges over partitions of $w$ into exactly $p$
  parts $\ge 1$, and $m_h(\lambda)$ counts the multiplicity of $h$ in
  $\lambda$.  Each mode level $h$ contributes $d$ generators; choosing
  $m_h$ of them and antisymmetrizing gives $\binom{d}{m_h}$.  The
  product over mode levels reflects the independence of distinct modes.

  For $d = 3$, the first few values are:
  $\dim \CE^1_w = 3$ for all $w \ge 1$ (three generators per mode);
  $\dim \CE^2_2 = 3$ (choosing two generators from the same mode $h = 1$:
  $\binom{3}{2} = 3$); $\dim \CE^2_3 = 9$ (partitions $1+2$:
  $\binom{3}{1}\binom{3}{1} = 9$); and so on.  At large $w$, the
  cochain dimensions grow as $\sim d^p \cdot p(w,p)/p!$ where $p(w,p)$
  is the number of partitions of $w$ into $p$ parts, which is polynomial
  in $w$ for fixed $p$.

  The concentration ratio $\rho_w$ divides the total cohomology (which
  depends on the bracket) by the total cochain dimension (which does not).
  For semisimple $\fg$ at a non-triangular weight $w$, the numerator
  is $0$ while the denominator grows polynomially, giving $\rho_w = 0$.
  The concentration is thus exponentially sparse in the weight direction:
  nonzero ratios occur only at the triangular weights $w = p(p+1)/2$,
  which grow quadratically in $p$.
\end{remark}


\subsection{$\mathfrak{sl}_3$: multi-orbit concentration}
\label{sec:sl3}

For $\fg = \mathfrak{sl}_3$ ($d = 8$), the affine Weyl group is the affine
$A_2$ group, with more than one minimal representative at higher lengths.
The computed bidegree data through weight $4$ is:

\begin{center}
  \begin{tabular}{c|ccccc}
    \toprule
    $p \backslash w$ & $1$ & $2$ & $3$ & $4$ \\
    \midrule
    $1$ & $8$ &  &  &  \\
    $2$ &  & $20$ &  &  \\
    $3$ &  &  &  & $63$ \\
    \bottomrule
  \end{tabular}
\end{center}

\noindent
At $p = 1$: weight $1$, dimension $8$ (the adjoint of $\mathfrak{sl}_3$).
At $p = 2$: weight $2$, dimension $20$.
At $p = 3$: weight $4$, dimension $63$.

The weight at $p = 3$ is $4$, not the triangular number $6$ that would
appear for $\mathfrak{sl}_2$ at the same CE degree.  This reflects the
$A_2$ root geometry: the affine Weyl group of $\mathfrak{sl}_3$ places the
length-$3$ contribution at a different weight than $A_1$ does.

The total dimensions are $8, 20, 0, 63, \ldots$\,, exhibiting vanishing at
weight $3$.  The growth rate is approximately $8^n$, reflecting the
exponential growth of $\dim \mathfrak{sl}_3$-modules indexed by the affine
Weyl group (contrast with the polynomial growth $2p + 1$ for
$\mathfrak{sl}_2$).

\begin{remark}\label{rem:euler-sl3}
  The Euler characteristic $\chi_w$ for $\mathfrak{sl}_3$ comes from
  $\prod_{n \ge 1}(1 - t^n)^8$.  At $w = 2$: $\chi_2 = \binom{8}{2} - 8 =
  28 - 8 = 20$, matching the single entry $H^2_2 = 20$.  At $w = 3$:
  $\chi_3 = -\binom{8}{3} + 8 \cdot 8 - 8 = -56 + 64 - 8 = 0$, consistent
  with the vanishing of total cohomology at weight $3$.
\end{remark}


\subsection{Tables of the differential rank}
\label{sec:tables}

\begin{table}[ht]
  \centering
  \begin{tabular}{cc|ccc}
    \toprule
    $p$ & $w$ & $\beta(\mathfrak{sl}_2)$ & $\beta(\fh_3)$ & $\beta(\kk^3)$ \\
    \midrule
    $1$ & $2$ & $1.000$ & $0.333$ & $0$ \\
    $1$ & $3$ & $1.000$ & $0.333$ & $0$ \\
    $1$ & $4$ & $1.000$ & $0.333$ & $0$ \\
    $2$ & $3$ & $0.111$ & $0$ & $0$ \\
    $2$ & $4$ & $0.750$ & $0.333$ & $0$ \\
    $2$ & $5$ & $0.833$ & $0.389$ & $0$ \\
    \bottomrule
  \end{tabular}
  \caption{Bracket rank fractions $\beta(p, w) = \rk(d^p_w)/\dim \CE^p_w$ for
  three-dimensional algebras.  For $\mathfrak{sl}_2$, $\beta(1, w) = 1$ at
  all $w \ge 2$: the first CE differential always has full rank.  For
  abelian algebras, $\beta \equiv 0$.}
  \label{tab:rank-fractions}
\end{table}


% ======================================================================
\section{The derived series and solvability spectrum}
\label{sec:derived}
% ======================================================================

The three mechanisms are detected by the algebraic invariants of $\fg$.  We
collect the structural data for the algebras under consideration.

\begin{table}[ht]
  \centering
  \begin{tabular}{lccccc}
    \toprule
    Algebra & $d$ & Derived series & LCS & $\pi$ & Class \\
    \midrule
    $\mathfrak{sl}_2$ & $3$ & $(3, 3)$ & $(3, 3)$ & $1$ & Semisimple \\
    $\mathfrak{sl}_3$ & $8$ & $(8, 8)$ & $(8, 8)$ & $1$ & Semisimple \\
    $\fb(\mathfrak{sl}_2)$ & $2$ & $(2, 1, 0)$ & $(2, 1, 1, \ldots)$ & $1/2$
    & Solvable \\
    $\fh_3$ & $3$ & $(3, 1, 0)$ & $(3, 1, 0)$ & $1/3$ & Nilpotent \\
    $\fn_+(\mathfrak{sl}_3)$ & $3$ & $(3, 1, 0)$ & $(3, 1, 0)$ & $1/3$ &
    Nilpotent \\
    $\kk^d$ & $d$ & $(d, 0)$ & $(d, 0)$ & $0$ & Abelian \\
    \bottomrule
  \end{tabular}
  \caption{Structural invariants of the test algebras.  The derived series
  is $(\dim \fg^{(0)}, \dim \fg^{(1)}, \ldots)$; the LCS is
  $(\dim \fg^0, \dim \fg^1, \ldots)$.  Semisimple algebras have constant
  derived series (perfectness); solvable algebras terminate at $0$;
  nilpotent algebras have LCS terminating at $0$.  The Borel $\fb$ is
  solvable but \emph{not} nilpotent: its LCS stabilizes at $1 \neq 0$.}
  \label{tab:derived}
\end{table}

\begin{remark}[The nilradical $\fn_+(\mathfrak{sl}_3)$]\label{rem:nilrad}
  The nilradical $\fn_+ = \operatorname{span}(e_1, e_2, e_{12})$ with
  $[e_1, e_2] = e_{12}$ is isomorphic to the Heisenberg algebra $\fh_3$ as
  a Lie algebra.  Accordingly, their loop CE cohomologies coincide weight by
  weight, and they have identical concentration ratios.  This confirms that
  the concentration depends on the \emph{abstract} Lie algebra structure,
  not on the ambient embedding.
\end{remark}


% ======================================================================
\section{Implications for chiral algebras}
\label{sec:chiral}
% ======================================================================

\subsection{Bar-cobar duality and Koszulness}

For the universal affine vertex algebra $V_k(\fg)$ at non-critical level
$k \neq -h^\vee$, the bar complex $B(V_k(\fg))$ is a factorization coalgebra
on the Ran space of a smooth curve $X$.  The PBW filtration on $V_k(\fg)$
induces a spectral sequence whose $E_1$ page is precisely the CE cohomology
$H^\bullet(\gminus, \kk)$ studied in this paper.

\begin{proposition}\label{prop:koszul-from-concentration}
  If $\fg$ is semisimple, then $V_k(\fg)$ is chirally Koszul at generic
  level $k$: the PBW spectral sequence degenerates at $E_2$, and the bar
  cohomology $H^\bullet(B(V_k(\fg)))$ is concentrated in bar degree~$1$.
\end{proposition}

\begin{proof}
  The Garland--Lepowsky concentration (Theorem~\ref{thm:GL}) shows that
  the $E_1$ page has a single nonzero row in each weight column at each
  CE degree.  The potential $E_2$ differentials connect different CE
  degrees, but the concentration forces them to map between weight spaces
  that are widely separated (the weight gaps grow linearly for
  $\mathfrak{sl}_2$ and at least linearly for higher rank).  The PBW
  degeneration at $E_2$ follows by the argument of~\cite[Proposition
  4.2.1]{L-MKD}.
\end{proof}

This proposition connects the ``why'' of Garland--Lepowsky concentration to
the structural foundation of chiral Koszul duality: the same mechanisms
(perfectness, Killing nondegeneracy, Weyl orbit uniqueness) that force CE
concentration also ensure that the chiral bar complex has the Koszul
regularity needed for the full bar-cobar duality theory.

The converse direction is equally informative: non-semisimple loop algebras
fail to concentrate (as shown in \S\ref{sec:counterexamples}), and
correspondingly the vertex algebras built from non-semisimple Lie algebras
(free bosons from abelian $\fg$, $\beta\gamma$ systems, lattice algebras)
have qualitatively different bar-complex behavior.  Free-field vertex
algebras are chirally Koszul but their bar cohomology concentrates for
reasons unrelated to the CE mechanism: the concentration comes from the
quadratic OPE structure rather than from perfectness of the underlying
Lie algebra.

\begin{remark}[Concentration and shadow depth]\label{rem:shadow-depth}
  The shadow depth $r_{\max}$ classifies the complexity of a chirally
  Koszul algebra: class~G (Gaussian, $r_{\max} = 2$) for free fields,
  class~L (Lie/tree, $r_{\max} = 3$) for affine algebras of semisimple
  $\fg$, class~C (contact, $r_{\max} = 4$) for $\beta\gamma$ systems,
  and class~M (mixed, $r_{\max} = \infty$) for the Virasoro and
  $W$-algebras.  Affine vertex algebras of semisimple $\fg$ belong to
  class~L precisely because the Garland--Lepowsky concentration forces
  the cubic shadow invariant $C(\fg)$ to be nonzero (the tree-level
  obstruction from the non-abelian bracket) but the quartic and higher
  shadows to vanish (the Killing form controls all higher obstructions
  through the root-lattice constraint).  Shadow depth $3$ is the
  algebraic fingerprint of the three-mechanism concentration.
\end{remark}


\subsection{The modular characteristic}

The modular characteristic of an affine vertex algebra is
\[
  \kappa(V_k(\fg)) = \frac{\dim(\fg) \cdot (k + h^\vee)}{2h^\vee},
\]
where $h^\vee$ is the dual Coxeter number.  This invariant controls the
genus expansion of the shadow obstruction tower: $F_g(V_k(\fg)) = \kappa
\cdot \lambda_g$ at all genera for uniform-weight algebras.  The
Garland--Lepowsky concentration ensures that $\kappa$ is computed from the
\emph{same} CE cohomological data that controls the spectral sequence
degeneration.  In this sense, the three mechanisms of
\S\ref{sec:mechanisms} underlie both the weight-level concentration and the
genus-tower universality of the shadow obstruction calculus.


\subsection{Chiral Koszul dual and generating functions}

For $\fg = \mathfrak{sl}_2$, the chiral Koszul dual is generated by bar
cohomology $H^1(B(V_k(\mathfrak{sl}_2)))$, whose weight-$n$ dimension is
the Riordan number $R(n+3)$ (OEIS A005043).  This sequence satisfies the
algebraic equation $x(1+x)R^2 - (1+x)R + 1 = 0$ and has the generating
function
\[
  P_{\mathfrak{sl}_2}(x) = \frac{1 + x - \sqrt{1 - 2x - 3x^2}}{2x(1+x)},
\]
which is algebraic of degree~$2$ over $\kk(x)$.  The discriminant
$\Delta(x) = 1 - 2x - 3x^2 = (1 - 3x)(1+x)$ is shared by all rank-$1$
standard chiral algebras (Virasoro, $\beta\gamma$), a universality that
extends the Garland--Lepowsky structure to the full chiral Koszul dual.

For $\fg = \mathfrak{sl}_3$, the bar cohomology generating function is
conjectured to be rational:
$P_{\mathfrak{sl}_3}(x) = 4x(2 - 13x - 2x^2)/((1-8x)(1-3x-x^2))$.
The exponential growth rate $\gamma \approx 8$ reflects the dimension
of $\mathfrak{sl}_3$ through the formula $\gamma = (d-1) + O(1)$, a
further manifestation of the root-system dependence revealed by
the Garland--Lepowsky analysis.

\begin{remark}[Distinction between CE and bar cohomology]\label{rem:CE-vs-bar}
  The CE cohomology $H^p(\gminus, \kk)_w$ and the bar cohomology
  $H^n(B(V_k(\fg)))_w$ are related but distinct objects.  The CE
  complex computes the cohomology of the \emph{loop algebra} $\gminus$
  (a Lie algebra cohomology), bigraded by CE degree $p$ and weight $w$.
  The bar complex computes the cohomology of the \emph{vertex algebra}
  $V_k(\fg)$ (a derived invariant of the chiral algebra), bigraded by
  bar degree $n$ and weight $w$.  For Koszul algebras, the bar
  cohomology is concentrated in bar degree $n = 1$, and the PBW
  spectral sequence relates the two: the $E_1$ page consists of the CE
  groups, and the $E_\infty$ page gives the bar groups.  At $E_2$
  degeneration (the Koszul case), the bar cohomology $H^1_w$ is
  determined by the CE data through the spectral sequence.  The Riordan
  numbers $R(n+3)$ giving $\dim H^1_n(B(V_k(\mathfrak{sl}_2)))$ encode
  the cumulative effect of CE concentration across all CE degrees at
  each weight; the Garland--Lepowsky dimensions $2p+1$ at triangular
  weights are the CE-level data that feeds into this determination.
\end{remark}


\subsection{The Euler invariant and its limitations}

The product formula~\eqref{eq:euler-product} is universal in $d = \dim\fg$
and blind to the bracket structure.  We can now state precisely what it
captures and what it misses.

\begin{proposition}\label{prop:euler-universal}
  Let $\fg$ and $\fg'$ be two Lie algebras of the same dimension $d$.  Then:
  \begin{enumerate}[label=(\alph*)]
    \item $\chi_w(\fg) = \chi_w(\fg')$ for all $w \ge 0$ (Euler universality).
    \item $\sum_p \dim H^p(\gminus, \kk)_w$ can differ between $\fg$ and
    $\fg'$ for $w \ge 2$ whenever $\fg$ and $\fg'$ have different bracket
    structures.
    \item For semisimple $\fg$, the Euler characteristic determines the
    individual dimensions: $|\chi_w| = \dim H^*(\gminus,\kk)_w$ at every
    weight $w$, because the CE cohomology is concentrated in a single
    CE degree at each weight.
  \end{enumerate}
\end{proposition}

\begin{proof}
  Part~(a) is the identity~\eqref{eq:euler-product}.  Part~(b) follows from
  the counterexamples: $\mathfrak{sl}_2$ and $\kk^3$ both have $d = 3$ and
  identical $\chi_w$, yet their total dimensions differ at $w = 2$ ($0$ vs
  $6$).  For part~(c), Theorem~\ref{thm:GL} states that at each triangular
  weight $w = p(p+1)/2$, only $H^p$ is nonzero, so $|\chi_w| = |(-1)^p
  \cdot (2p+1)| = 2p+1 = \dim H^p_w$.  At non-triangular weights, $\chi_w =
  0$ and all $H^p_w = 0$.
\end{proof}

Part~(c) is the ``accidental'' equality alluded to in the introduction:
for semisimple $\fg$, the absolute value of the Euler coefficient at each
weight equals the total cohomology dimension.  This coincidence fails for
every non-semisimple algebra of the same dimension, because the lack of
concentration allows positive and negative contributions to partially cancel
in the Euler sum while leaving large total dimensions.


% ======================================================================
\section{Proof of the main theorem}
\label{sec:proof}
% ======================================================================

We consolidate the results into a single statement.

\begin{theorem}[Three mechanisms of concentration]\label{thm:main}
  Let $\fg$ be a finite-dimensional Lie algebra of dimension $d$ over a
  field of characteristic zero.  The following are equivalent:
  \begin{enumerate}[label=(\roman*)]
    \item $\fg$ is semisimple.
    \item The bracket surjectivity index $\sigma_w(\fg) = 1$ for all $w \ge 2$.
    \item The Killing form $B(x,y) = \tr(\ad_x \circ \ad_y)$ is nondegenerate.
    \item The CE cohomology $H^\bullet(\gminus, \kk)$ is concentrated in the
    sense that for each CE degree $p$, the number of weights $w$ with
    $H^p(\gminus, \kk)_w \neq 0$ equals $|\Waff^p|$ (the number of
    minimal affine Weyl coset representatives of length~$p$).
  \end{enumerate}
  For algebras of the same dimension $d$, the concentration
  ratio satisfies
  \[
    \rho_w(\fg_{\mathrm{ab}}) \ge \rho_w(\fg_{\mathrm{nil}})
    \ge \rho_w(\fg_{\mathrm{ss}})
  \]
  at every weight $w$, where $\fg_{\mathrm{ab}}$ is abelian,
  $\fg_{\mathrm{nil}}$ is nilpotent with nontrivial bracket, and
  $\fg_{\mathrm{ss}}$ is semisimple.
\end{theorem}

\begin{proof}
  The equivalence (i)$\Leftrightarrow$(ii) follows from the fact that
  $\sigma_w = \pi(\fg)$ for all $w \ge 2$ (as observed in
  \S\ref{sec:perfectness}), and $\pi(\fg) = 1$ if and only if $\fg$ is
  perfect.  For finite-dimensional Lie algebras in characteristic zero,
  perfectness is equivalent to semisimplicity (since the radical of a
  non-semisimple algebra is a proper ideal mapping to zero under the
  bracket projection).

  The equivalence (i)$\Leftrightarrow$(iii) is Cartan's criterion.

  The implication (i)$\Rightarrow$(iv) is the Garland--Lepowsky theorem
  (Theorem~\ref{thm:GL}).  For (iv)$\Rightarrow$(i): if $\fg$ is not
  semisimple, then the Killing form is degenerate, and Propositions
  \ref{prop:abelian-no-concentration} and~\ref{prop:borel-test}
  show that non-semisimple algebras fail to achieve the Weyl-orbit
  concentration pattern.  Specifically, the Borel test
  (\S\ref{sec:borel}) demonstrates that even solvable algebras with
  nontrivial brackets have cohomology at weights outside the Weyl orbit
  loci.

  The concentration ordering follows from Theorem~\ref{thm:ordering}.
\end{proof}

\begin{remark}\label{rem:three-necessary}
  The three mechanisms are independently necessary in the following sense:
  \begin{enumerate}[label=(\alph*)]
    \item The Heisenberg algebra has a nontrivial bracket ($\sigma = 1/3$)
    but degenerate Killing form and no Weyl group.  Mechanism~(I) alone is
    insufficient.
    \item The Borel subalgebra has $\sigma = 1/2$ (stronger than Heisenberg)
    but degenerate Killing form and no Weyl group.  Mechanism~(I) at
    partial strength is still insufficient.
    \item No non-semisimple algebra possesses a nondegenerate Killing form
    (Cartan's criterion), so mechanisms~(II) and~(III) are available only
    for semisimple algebras.
  \end{enumerate}
  All three mechanisms engage simultaneously if and only if $\fg$ is
  semisimple.
\end{remark}


% ======================================================================
\section{Computational verification}
\label{sec:verification}
% ======================================================================

All computations in this paper are independently verified by the
\texttt{garland\_lepowsky\_why\_engine} module (77~tests) and the
\texttt{bar\_loop\_group\_engine} module (62~tests) in the repository
\texttt{compute/}.  The verification follows the multi-path protocol:
every numerical value is confirmed by at least two independent computation
methods.

\begin{enumerate}[label=(\arabic*)]
  \item The Garland--Lepowsky formula~\eqref{eq:GL-sl2} is verified by
  direct CE matrix computation at weights $w \le 10$ and independently by
  the Euler product cross-check (signed sums match $\prod(1-t^n)^3$).

  \item The concentration ratios in Table~\ref{tab:concentration-d3} are
  computed from CE cohomology dimensions (path~1) and independently from
  the ratio of total cohomology to total cochain dimensions (path~2).

  \item The Killing form determinants are computed by explicit matrix
  construction ($\ad$-matrices and trace products) and independently verified
  by the rank computation.

  \item The Borel test (Proposition~\ref{prop:borel-test}) is verified by
  comparing Borel and abelian total dimensions weight by weight, with the
  Euler series cross-check confirming that signed sums agree.

  \item The bracket surjectivity indices are computed from the rank of
  the bracket image matrix and independently verified to equal the
  perfectness ratio.

  \item The $\mathfrak{sl}_3$ bidegree data is verified by direct $8
  \times 8$ CE matrix computation and independently by the Euler series
  cross-check ($\prod(1-t^n)^8$ coefficients).
\end{enumerate}


% ======================================================================
\section{Open questions}
\label{sec:open}
% ======================================================================

The analysis in this paper suggests several directions for further
investigation.

\begin{enumerate}[label=(\arabic*),leftmargin=*]

  \item \textbf{Quantitative concentration for higher-rank algebras.}
  The $\mathfrak{sl}_2$ case has a single Weyl orbit per CE degree
  ($|\Waff^p| = 1$), producing the cleanest concentration.  For
  $\mathfrak{sl}_N$ with $N \ge 3$, the affine Weyl group $\Waff(A_{N-1})$
  has $|\Waff^p| > 1$ at higher lengths, so multiple weights contribute at
  each CE degree.  The computed data for $\mathfrak{sl}_3$ shows
  concentration that is still strong (zero total cohomology at $w = 3$) but
  less sharp than $\mathfrak{sl}_2$.  A precise formula for
  $|\Waff^p(A_{N-1})|$ as a function of $p$ and $N$ would quantify how
  concentration degrades with rank.  For exceptional types, the orbit
  structure is more complex, and the concentration pattern has not been
  computed beyond the first few weights.

  \item \textbf{Concentration at critical level.}
  The Garland--Lepowsky theorem applies to the loop algebra $\gminus$
  independently of the level $k$ of the affine vertex algebra $V_k(\fg)$.
  At the critical level $k = -h^\vee$, the vertex algebra develops a large
  center (the Feigin--Frenkel center $\mathfrak{z}(\hat{\fg})$), and the PBW
  spectral sequence may have different degeneration behavior.  The CE
  concentration persists (it depends only on $\fg$, not $k$), but the
  passage from CE to bar cohomology may involve higher differentials.
  Whether bar concentration holds at critical level is an open question
  with implications for the geometric Langlands programme.

  \item \textbf{Refinement of the concentration ordering.}
  Theorem~\ref{thm:ordering} establishes $\rho(\text{abelian}) \ge
  \rho(\text{nilpotent}) \ge \rho(\text{semisimple})$.  Can this be
  refined to a continuous ordering indexed by the perfectness ratio $\pi$?
  For algebras of the same dimension, is $\rho_w(\fg)$ a monotone
  decreasing function of $\pi(\fg)$ at each weight $w$?  The computed
  examples support this, but the Borel and Heisenberg algebras have
  different dimensions ($2$ and $3$), so a direct comparison requires
  normalization.

  \item \textbf{Characteristic $p$ phenomena.}
  The entire analysis assumes characteristic zero.  In positive
  characteristic, Weyl's complete reducibility theorem fails, the Killing
  form may be degenerate for simple algebras (the first Whitehead lemma
  fails), and the Garland--Lepowsky concentration may break down.  The
  $\mathfrak{sl}_2$ concentration at weight $p(p+1)/2$ with dimension
  $2p+1$ suggests a connection to the representation theory of $SL_2$
  in characteristic $p$ (where $V_p$ becomes reducible).  An investigation
  of CE cohomology concentration in positive characteristic would
  illuminate the role of complete reducibility in the mechanism.

  \item \textbf{Infinite-dimensional generalizations.}
  The Virasoro algebra is perfect ($[L_m, L_n] = (m-n)L_{m+n} + \ldots$
  is surjective onto $\text{Vir}/\CC c$) and has a nondegenerate Shapovalov
  form (the analogue of the Killing form).  Yet it is not semisimple in the
  classical sense.  Whether the Virasoro loop algebra exhibits
  Garland--Lepowsky-type concentration is an interesting question; the
  shadow depth $r_{\max} = \infty$ for the Virasoro vertex algebra
  suggests that the concentration pattern is qualitatively different.

\end{enumerate}


\begin{thebibliography}{10}

\bibitem{GL76}
H.~Garland and J.~Lepowsky,
\emph{Lie algebra homology and the Macdonald--Kac formulas},
Invent.\ Math.\ \textbf{34} (1976), 37--76.

\bibitem{Kostant61}
B.~Kostant,
\emph{Lie algebra cohomology and the generalized Borel--Weil theorem},
Ann.\ of Math.\ (2) \textbf{74} (1961), 329--387.

\bibitem{Kumar02}
S.~Kumar,
\emph{Kac--Moody Groups, their Flag Varieties and Representation Theory},
Progress in Mathematics, vol.\ 204, Birkh\"auser, 2002.

\bibitem{BottSamelson58}
R.~Bott and H.~Samelson,
\emph{Applications of the theory of Morse to symmetric spaces},
Amer.\ J.\ Math.\ \textbf{80} (1958), 964--1029.

\bibitem{L-MKD}
R.~Lorgat,
\emph{Modular Koszul Duality},
in preparation, 2026.

\bibitem{FBZ04}
E.~Frenkel and D.~Ben-Zvi,
\emph{Vertex Algebras and Algebraic Curves}, 2nd ed.,
Mathematical Surveys and Monographs, vol.\ 88, AMS, 2004.

\bibitem{Weyl25}
H.~Weyl,
\emph{Theorie der Darstellung kontinuierlicher halbeinfacher Gruppen durch
lineare Transformationen},
Math.\ Z.\ \textbf{23} (1925), 271--309; \textbf{24} (1926), 328--376,
377--395.

\bibitem{Kac90}
V.~G.~Kac,
\emph{Infinite-dimensional Lie Algebras}, 3rd ed.,
Cambridge University Press, 1990.

\end{thebibliography}

\end{document}
